\chapter{DNA 的结构为什么适合保存信息}

\begin{kaichang}
厨房里就能完成一个听起来像科幻的实验。取一颗熟透的草莓，装进保鲜袋捏碎，加一点儿盐水、几滴洗洁精，轻轻晃一晃，再用纱布滤出汁液。然后，把冰过的酒精缓缓倒在滤液表面，静置几分钟——液面交界处会慢慢浮起一团白色絮状物，像一小团棉线。拿筷子尖去挑，它居然能被缠起来，越绕越多。

先猜一猜：这团白线是什么？蛋白质？糖？果肉纤维？

都不是。那是草莓的 DNA——从几十亿个细胞里放出来的遗传分子。你刚刚用筷子，把生命的“说明书”钓出了水面。这一章要回答的问题是：这种分子到底长什么样，为什么它的这副长相，恰好适合干“保存信息”这份工作？
\end{kaichang}

\section{先从草莓里钓出一团白线}

先只描述眼睛和手能感觉到的东西。那团白色絮状物黏糊糊的，能拉丝，绕在筷子上像一小段透明渔线。它溶于水却不溶于冷酒精，所以会在酒精和汁液的交界处“现身”——交界处先出现一层薄雾，雾慢慢聚拢成絮。滤出的汁液本身是粉红色的，但絮状物是白的，说明它既不是色素也不是果肉。草莓越熟、捏得越碎，得到的絮状物越多。选草莓还有一个隐秘的理由：栽培草莓是八倍体，每个细胞里有八套染色体（第 64 章讲过染色体是遗传物质的“集装箱”），DNA 含量比一般水果高得多，所以特别适合做这个实验。

\begin{shiyan}
\textbf{材料}：熟草莓两三颗、食盐、水、洗洁精、纱布或咖啡滤纸、透明玻璃杯、医用酒精（提前放冰箱冷冻室冰镇，不要冻成冰）、筷子一根。

\textbf{步骤}：① 草莓去蒂装袋，捏成糊状；② 半杯温水化开一小勺盐，加几滴洗洁精，倒入袋中轻晃两分钟（不要猛摇出泡）；③ 用纱布过滤到玻璃杯里；④ 沿杯壁缓缓倒入等体积的冰酒精，静置 5 分钟；⑤ 用筷子尖在交界处轻轻搅动，把絮状物缠上来。

\textbf{观察点}：絮状物出现的速度、能否拉丝、在酒精里和在水里的溶解差别。想一想：为什么要加盐？为什么要冰酒精？

\textbf{安全提示}：酒精易燃，全程远离明火；实验用过的材料不要再吃；洗洁精别弄到眼睛里。这个实验在家可以做，但做完要洗手。
\end{shiyan}

肉眼看不见的东西是：那一根“线”并不是一个分子，而是数以亿计的超长分子纠缠在一起。单个 DNA 分子细得只有约 2 纳米，比可见光的波长还小几百倍，任何光学显微镜都看不清它的细节。同样的方法对香蕉、洋葱、豌豆都管用——凡是细胞多、DNA 多的材料，都能钓出白线，只是草莓特别大方。要看清这副分子的长相，得先把镜头推进到原子层面。

\section{把镜头推进：一块核苷酸}

DNA 是一种高分子，也就是由许多重复单元连成的大分子（第 43 章的塑料、第 46 章的淀粉都是高分子）。它的重复单元叫\textbf{核苷酸}。每个核苷酸由三部分拼成，像一枚三合一的积木：

\begin{itemize}[itemsep=2pt]
\item 一个\textbf{磷酸}：就是第 14 章认识的磷的含氧酸根，在水中带负电；
\item 一个\textbf{脱氧核糖}：五个碳原子连成的小环，属于糖（第 46 章），只是比核糖少了一个氧原子，所以叫“脱氧”；
\item 一个\textbf{碱基}：含氮的扁平环状分子（第 14 章的氮、第 41 章的芳香环在这里汇合），是积木上真正“写字”的部分。
\end{itemize}

核苷酸这种积木，你其实早就“吃过”也“用过”。第 50 章的能量货币 ATP，就是一个腺嘌呤核苷酸背上三个磷酸；让香菇、鸡汤鲜美的呈味核苷酸，也是这类分子的亲戚——核苷酸不是实验室里的稀罕物，而是生命最通用的零件之一。“核酸”这个名字则来自历史：它最早从细胞核里被分离出来，而且呈酸性（骨架上那些磷酸的功劳）。

碱基只有四种：腺嘌呤（A）、鸟嘌呤（G）、胞嘧啶（C）、胸腺嘧啶（T）。A 和 G 是双环结构，个头大，叫嘌呤；C 和 T 是单环，个头小，叫嘧啶。四种字母，就是 DNA 全部的文字。

核苷酸们怎样连成链？磷酸像一截一截的连接件，把一个脱氧核糖和下一个脱氧核糖串起来，于是链的“龙骨”是糖—磷酸、糖—磷酸交替排列的骨架，四种碱基则像旗子一样从骨架的每个糖上伸出来。这条骨架有两副脾气要记住：磷酸带负电，所以整条骨架亲水、爱跟正离子打交道；而伸出去的碱基是扁平的环，比较“怕水”（回想第 47 章脂质分子的疏水脾气）。这两副脾气，待会儿会决定两条链怎样拥抱在一起。

\section{规则一：链有方向，配对有规矩}

先看方向。脱氧核糖环上的五个碳原子按惯例编了号，从 $1'$ 到 $5'$（读作“一撇”到“五撇”）。磷酸总是把这个糖的 $3'$ 号碳和下一个糖的 $5'$ 号碳连起来，一节一节走下去。于是整条链有了头和尾：一端挂着自由的 $5'$ 磷酸，叫 $5'$ 端；另一端露出自由的 $3'$，叫 $3'$ 端。DNA 链就像一列火车，车厢有前后，列车有车头车尾——后面会看到，复制和读取的机器都要先认清方向才能开工。

再看配对。两条 DNA 链并排时，碱基两两相对，而且搭配极其严格：A 只和 T 配对，中间形成两个氢键；G 只和 C 配对，中间形成三个氢键（氢键是第 22、23 章讲过的分子间弱相互作用）。为什么是这两对？因为它们的氢键“供体”和“受体”位置恰好互补，就像插头和插座：A 伸出氢原子的位置，正对 T 上孤对电子的位置，换一对就对不上。还有一个几何上的讲究：嘌呤大、嘧啶小，一大一小搭配，每一对的宽度几乎一样，整条双链才能保持约 2 纳米的均匀粗细，不会忽粗忽细。

最后看两条链的相对方向：它们是\textbf{反向平行}的——一条从 $5'$ 到 $3'$ 走，另一条就从 $3'$ 到 $5'$ 走，像两股方向相反的扶梯。这个细节看似琐碎，却是第 67 章复制机器的“命门”，先记住它。

第 65 章讲过的查加夫规则，在这里可以再看一眼：任何来源的双链 DNA 里，A 的数量总等于 T，G 总等于 C。当时这只是一串奇怪的统计数字，没有解释；放进配对模型里，它立刻变成必然——每一个 A 都牵着对面一个 T。一个好理论的标志就是这样：它不只解释新现象，还顺手把旧的悬案变成“本该如此”。

两条链不是并排躺着就算了。配对后，整条双链会自然扭转成螺旋，每转一圈大约 10 个碱基对——这就是第 65 章讲的、由 X 射线衍射照片揭示的\textbf{双螺旋}。那章回答了“科学家怎么知道是双螺旋”；这一章要回答的是：这副结构凭什么适合保存信息。

\section{规则二：弱，才是设计的精髓}

很多人第一反应是：保存信息当然要越牢固越好，两条链之间应该“粘”得死死的。恰恰相反。DNA 的高明之处，在于把“弱”和“多”组合到了极致。

\begin{qiaoliang}
来比一比能量尺度。单个氢键的结合能大约只有 \ud{10\sim 20}{kJ\,mol^{-1}}，而一根共价键大约 \ud{350}{kJ\,mol^{-1}}，体温下分子的平均热运动能量也有约 \ud{2.5}{kJ\,mol^{-1}} 的数量级。也就是说，三五个氢键握在一起，随时可能被热运动摇开——实际上，DNA 双螺旋在体温下一直在局部地“呼吸”：小段的配对不断打开又合上。但一段 1000 个碱基对的 DNA 有两三千个氢键，要同时摇开它们，概率小到可以忽略。这就像拉链：一个齿轻轻一拨就开，整条拉链却很牢靠；而你想拉开它时，从头开始逐个解开又毫不费力。碱基之间除了氢键，还有一层更沉默的贡献：扁平的碱基像一摞硬币一样叠起来，环与环之间的范德华力加上“怕水”的碱基躲进螺旋内部、避开水分子所获得的稳定（第 23、47 章的疏水效应），才是双螺旋整体稳定的主要来源。氢键负责“认人”——保证 A 只找 T；堆积负责“抱团”——让整条链不散架。一个管特异性，一个管稳定性，分工明确。
\end{qiaoliang}

现在可以回答章标题的一半了。DNA 结构适合保存信息，至少有三条结构理由：

\begin{itemize}[itemsep=2pt]
\item \textbf{字母少而序列自由}。骨架千篇一律，四种碱基的排列却不受限制，任何“文章”都写得出来。
\item \textbf{每条链都是完整备份}。配对规则意味着：知道一条链的序列，就能唯一推出另一条。信息存了两份，而且两份可以拆开。
\item \textbf{弱键集合允许“读写”}。需要复制或修复时，分子机器可以像拉拉链一样局部打开双链，以一条链为模板拼出它的互补链；用完再合上。存得稳和取得出，这两个看似矛盾的要求，被“大量弱相互作用”同时满足了。
\end{itemize}

第二条还顺便解决了“坏了怎么办”：如果一条链上的某个碱基受损，细胞可以拿完好的另一条链当标准答案来对照修复。损坏与修复是下一章的主题之一，这里只记住结构上的伏笔：互补双链天生自带纠错底稿。

\section{把双螺旋写成一行字母}

铺垫了这么多，符号终于可以登场了。因为骨架永远一样、配对永远固定，生物学家干脆只写一条链的碱基顺序，并标明方向，例如：

\[5'\text{-ATGCGTAC-}3'\]

与它配对的另一条链不用猜，直接按规则写出来，注意方向相反：

\[3'\text{-TACGCATG-}5'\]

把双螺旋画成“梯子”的示意图，是这样的（这是人为的表示法：真实的碱基不是字母，分子也没有颜色）：

\begin{center}
\begin{tikzpicture}[scale=0.9]
\draw[thick] (0,0) -- (0,4.4);
\draw[thick] (2.6,0) -- (2.6,4.4);
\foreach \y/\a/\b in {0.4/T/A, 1.2/C/G, 2.0/G/C, 2.8/A/T, 3.6/T/A}{
  \draw (0,\y) -- (2.6,\y);
  \node at (0.75,\y) {\small \a};
  \node at (1.85,\y) {\small \b};
}
\node at (-0.3,4.6) {\small $3'$};
\node at (-0.3,-0.2) {\small $5'$};
\node at (2.9,4.6) {\small $5'$};
\node at (2.9,-0.2) {\small $3'$};
\node at (1.3,-0.7) {\small 骨架（糖—磷酸）\quad 横档（碱基对，中间是氢键）};
\end{tikzpicture}
\end{center}

注意图中两条竖线的方向标注相反——反向平行就在这一笔里。真实分子会再拧成螺旋，但信息全在横档的顺序里，螺旋只是这架梯子最省空间、最稳当的收纳姿态。

一行字母能装多少信息？每个位置有 4 种选择，$n$ 个位置就有 $4^n$ 种序列。10 个位置已经超过一百万种（$4^{10}\approx 1.05\times 10^{6}$）；人类基因组的单套约 $3\times 10^{9}$ 个碱基对，可能的序列数是 $4$ 的三十亿次方，这个数字写出来需要的纸比整个可观测宇宙还多。换个更直观的账：把人类单套基因组的三十亿个字母印成书，一页排 1000 个字母，要印 300 万页，按每本 500 页装订，是 6000 本书——你身体几乎每个细胞里都揣着这样一座图书馆，而且是两套（第 64 章的双倍体）。

这些字母同时还占着物理空间。双螺旋里相邻碱基对相距 0.34 纳米，单套基因组拉直了长约 $3\times 10^{9}\times 0.34$ 纳米，约 1 米；两套约 2 米。一颗直径只有约 6 微米的细胞核，要塞下 2 米长的分子——这笔账怎么算，本章最后再收。

序列即信息，这条原理早已走出实验室，变成你新闻里常见的技术。亲子鉴定比对的是若干段因人而异的短序列：孩子的每一段都必然能在父母一方找到来源（第 63 章的概率规律在这里变成法庭证据）；刑侦里的 DNA 指纹同理，现场一根头发、一点唾液里的分子，就能和数据库里的序列对号。核酸检测则是另一条思路：人工合成一小段单链 DNA 当“探针”，利用的正是“A 只认 T、G 只认 C”的专一性——样本里若有病毒的互补序列，探针就会贴上去发出信号。这些技术没有一项需要看清分子的样子，全都只靠两件事：序列是信息，配对认死理。

\section{加热一下，证据就自己冒出来了}

\begin{zhengju}
双螺旋模型画得很美，但“碱基靠弱键配对、还能可逆分开”这件事，怎么验证？答案藏在一个便宜得多的实验里：给 DNA 水溶液慢慢加热，同时测量它对紫外线的吸收。

碱基会吸收波长 260 纳米附近的紫外线，但有个微妙之处：当它们整齐地叠在双螺旋内部时，吸收反而变弱；一旦双链分开、碱基暴露，吸收就明显增强，幅度可达约四成。20 世纪 50 年代末，研究者（如马默和多蒂等人）利用这个效应画出了 DNA 的“熔化曲线”：温度升高时，起初吸收几乎不动，到某个温度附近突然陡升——像冰在熔点化开一样，双链在这个温度范围内集中地解开成单链。这个温度被叫做熔链温度 $T_m$。

第一个关键发现：不同来源的 DNA，$T_m$ 不同，而且 G 和 C 含量越高的 DNA 越“耐热”，两者几乎成一条直线。为什么这能当证据？如果双链靠某种说不清的力量粘着，没理由偏偏跟 G、C 的比例挂钩；而按配对模型，G—C 对有三个氢键、堆积也更稳，自然更抗热——预言和测量严丝合缝。

第二个关键发现更漂亮：把分开的 DNA 慢慢降温，两条链会重新找到彼此，恢复双螺旋，紫外线吸收也降回去。这说明解开时没有弄坏任何“零件”，每条链都完整保存着认出对方所需的全部信息。更有说服力的是“杂交”实验：把来自两个不同物种的细菌的 DNA 混在一起解开再降温，亲缘近的物种之间会搭出不少杂合双链，亲缘远的几乎搭不起来——配对认的不是“随便一条链”，而是“序列互补的那一条”。这个方法后来甚至被用来给细菌“排家谱”。

也曾有别的解释被摆上台面：吸收变化会不会是蛋白质杂质在捣乱？但高度纯化的 DNA 照样出现这条曲线；会不会是加热把共价键烧断了？可那样就无法可逆复原。一个个替代解释被排除，留下的图景就是：特异性来自氢键配对，整体稳定还叠加了碱基堆积，双链可以可逆地开与合。今天你做的每一次核酸检测、每一轮 PCR，都在重复这个实验：升温到九十度出头让双链分开，降温让互补链对上——七十年前的熔化曲线，就是这些技术的说明书。
\end{zhengju}

\section{这个模型什么时候会失灵}

“字母写在两条互补链上”的图景很好用，但别把它当成 DNA 在细胞里的全部真相。它至少有四处会失灵：

\begin{itemize}[itemsep=2pt]
\item \textbf{不是所有遗传分子都是双链 DNA}。有些病毒的基因组是单链 DNA，还有一大批病毒干脆用 RNA（第 65 章已提醒过）。双螺旋是生命的“主流格式”，不是唯一格式。
\item \textbf{细胞里的 DNA 从不裸露}。教科书里那根光滑的梯子，在细胞中始终缠着蛋白质、打着超螺旋、被折叠压缩。模型描述的是局部结构，不是它在细胞里的日常姿态。
\item \textbf{字母会被磨损}。水分子的撞击、紫外线、活性氧，每天都在每个细胞的 DNA 上造成数以万计的损伤，比如胞嘧啶悄悄脱掉一个氨基变成别的碱基。如果没有一刻不停的修复（第 67 章细讲），这座图书馆几周就会变得错字连篇。“完美静态档案”的图景是错的——它更像一份被无数校对员日夜维护的手稿。
\item \textbf{双螺旋不止一种拧法}。常见的是右旋的 B 型；脱水时会出现更矮胖的 A 型，某些序列还能拧成左旋的 Z 型。结构会随环境和序列微调，不是一根僵硬的杆。
\end{itemize}

\begin{wuqu}
“两条链靠氢键粘在一起，所以氢键一定很强、很难分开。”——错在两个地方。第一，单个氢键在化学键家族里是出了名的弱，体温的热运动就能摇开一小段；双螺旋的稳定靠的是几千万个弱相互作用累加，再加上碱基堆积这个沉默的主力。第二，“容易分开”恰恰是功能的一部分：复制和读取都得先把链打开，如果两条链真的粘得死死的（比如全部换成共价键），信息倒是“保险”了，却永远取不出来——那不是保险柜，是棺材。好的信息存储介质，必须同时满足“平时散不了”和“用时打得开”，DNA 用“大量弱键”这个方案同时做到了两点。
\end{wuqu}

\section{再看筷子上的那团白线}

现在可以一步步解释草莓实验了，每种试剂都有它的岗位。洗洁精负责拆房子：它的分子一头亲水一头亲油，专门瓦解细胞膜和核膜的脂质双层（第 47、51 章），把细胞内容物放出来。盐负责劝架：DNA 骨架上密布的磷酸负电让分子互相排斥，盐里的正离子（第 18 章）围在骨架周围屏蔽电荷，分子们才肯靠近、聚到一起。冰酒精负责收网：DNA 靠骨架的亲水性溶在水里，酒精挤进来抢走水分子（第 25 章讲过“相似相溶”的另一面），DNA 的溶解度骤降，只好析出；低温让析出更彻底，也减慢捣乱酶的活动。最后你挑起的白线，是亿万个超长的 DNA 分子缠成的线团——单个分子细得看不见，亿万根缠在一起就成了肉眼可见的絮。

至于字母，你在筷子上一个也读不出来。这不奇怪：一个碱基只有纳米大小，那双“读出字母”的眼睛要到第 68、69 章讲细胞自己的读取机器、以及测序技术时才会配上。但你至少亲眼确认了第 65 章那个结论的物质形态：遗传信息真的是一种长长的、可以摸到的分子。顺便说一句：同样的白线也能从猪肝、豌豆苗甚至你自己漱口脱落的口腔细胞里钓出来——这个实验不分物种，因为所有细胞生命用的是同一套分子、同一副双螺旋。这本身就是“生命之树同根”的一件物证，等讲到进化（第十一篇）时我们还会回来用它。

\section{两米长的线，怎样装进微米大的核}

还剩那笔没算完的账。每个碱基对占 0.34 纳米，人类单套基因组拉直约 1 米，双倍体约 2 米，而细胞核直径只有约 6 微米——直线尺度上压缩了几十万倍。更夸张的总账：你的身体大约有 $10^{13}$ 个细胞，把所有 DNA 首尾相接，长约 $2\times 10^{13}$ 米，也就是约 200 亿公里，够在地球和太阳之间往返六十多趟。

细胞的解决方案是“分级缠线”，第 64 章见过的染色体就是它的成品。第一级：DNA 在一种叫组蛋白的蛋白质线轴上绕圈——组蛋白带正电，恰好抱住带负电的骨架，绕好的一节叫核小体，整串看起来像一条“珠链”，直径约 11 纳米。第二级：珠链再盘绕折叠成更粗的纤维。第三级：纤维打环、再层层压缩，到细胞分裂时折叠成显微镜下可见的致密染色体。就像把 2 米长的耳机线先缠在线轴上、再卷成小团塞进收纳盒。精子细胞把这套功夫推到了极限：它用比组蛋白更小的鱼精蛋白把 DNA 压成几乎致密的“行李卷”，因为精子只要送货，不打算在路上读书。

但请注意一个关键差别：收纳盒里的耳机线是死的，细胞核里的 DNA 是活的。哪一段要被读取，哪一段的包装就临时松开；哪一段长期不用，就压得更紧。包装的松紧本身就是一套调控开关——这是第 70 章“基因不是一直开着的灯”要展开的故事。

而压缩之前，还有一个更紧迫的工程：细胞每分裂一次，就得把这两米长、三十亿对字母的分子从头到尾抄写一遍，而且几乎不能抄错。凭什么能抄？靠的就是本章的主角——反向平行加互补配对：把两条链拉开，每一条都是完整的模板。具体由哪些分子机器动手、前头提到的那股“反向扶梯”会惹出什么麻烦、抄错了谁来校对——第 67 章，我们去看这场每秒数十亿次的书写。

\begin{tiaozhan}
实验室里常要估算一小段 DNA 的熔链温度，比如给 PCR 设计引物。对十几个字母的短链，有一条粗糙但好用的经验式：
\[T_m \approx 4\times (G+C\text{ 的个数}) + 2\times (A+T\text{ 的个数}) \quad (^{\circ}\mathrm{C})\]
它背后的思想就是本章的规则：G—C 对多一个氢键、堆积更稳，贡献约 4 度，A—T 对约 2 度。拿 $5'$-ATGCGT-$3'$ 试试：G、C 各 2 个共 4 个，A、T 各 1 个共 2 个，$T_m\approx 4\times 4+2\times 2=20\,^{\circ}\mathrm{C}$——这么短的链室温下就散了，难怪真正的引物通常要设计到 20 个字母上下。请自己再算一条 G、C 占一半的 20 聚短链，你会发现 $T_m$ 恰好落在六十度上下，这正是 PCR 降温配对步骤常用的温度区间。这条公式对长链不适用（长链要考虑盐浓度和链间竞争），但它干净地演示了：宏观的“熔点”，是微观每一个弱键的集体投票。跳过本框不影响主线。
\end{tiaozhan}

\section*{练一练}

\begin{sikaoti}
\item 写出与 $5'$-GATCCA-$3'$ 配对的另一条链的序列，并正确标注方向；再数一数这段双链里一共有多少个氢键。
\item 画一幅粒子示意图：两个核苷酸连成一小段骨架，标出磷酸、脱氧核糖、碱基的位置和 $5'$、$3'$ 端。图旁注明：圆圈和连线是人为表示法，不是分子的真实照片。
\item 为什么 G、C 含量高的 DNA 更耐热？请用“能量”和“概率”各解释一遍，不允许只说“因为 G—C 更牢固”。
\item 一段 1000 个碱基对的 DNA 拉直有多长（用纳米表示）？如果把它和直径约 10 微米的细胞并排比较，相当于多长的线放进多大的球？再想想：细胞核靠什么办法解决这个收纳问题？
\item 画一张信息流图：在草莓实验里，信息从染色体的碱基序列出发，到你筷子上的白线，经历了哪些“形态变化”？在哪一步之后，你肉眼再也“读”不到信息了？信息本身消失了吗？
\item 有人提议：“把两条链之间的氢键全部升级成共价键，遗传物质就更稳定了。”从“存得稳”和“取得出”两个要求出发，论证这个“升级”为什么是一场灾难。
\end{sikaoti}
